\documentclass[aps,prb,twocolumn,superscriptaddress,floatfix]{revtex4-2}
\usepackage{amsmath,amssymb,graphicx,physics,bm}
\begin{document}

\title{Replication Report: Light-driven octupolar inverse Faraday effect\\
and multipolar order in Mott insulators (Banerjee \emph{et al.}, 2026)}
\author{TEXTURES-100 Replication Subagent}
\affiliation{Automated from-scratch replication; kernel provenance:
\texttt{ollie\_multipolar\_stevens\_landau\_kernel.py}}
\date{\today}

\begin{abstract}
We independently reconstruct the central physics of Banerjee \emph{et al.}
(arXiv:2605.08049v1): circularly polarized light (CPL) applied to a
spin--orbit-coupled $4d^2/5d^2$ Mott insulator with edge-sharing octahedra
generates, via a Floquet Schrieffer--Wolff (van~Vleck high-frequency) expansion
of a driven Hubbard--Kanamori model, an \emph{effective static field} $h_m$ that
couples \emph{linearly} to the magnetic octupole $T_{xyz}$ (pseudospin
$\tilde\sigma^y$) --- the octupolar inverse Faraday effect (OIFE) --- together
with a bond-dependent anisotropic exchange $\Gamma^{(3)}$. We verify (i) the
$E_g$ non-Kramers pseudospin SU(2) algebra from the normalized Stevens
operators, (ii) the analytic Bessel--Floquet couplings $J_\mathrm{eff},
\Gamma^{(3)}, h_m$ (Eqs.~6a--6f), (iii) the van~Vleck mechanism
$H_\mathrm{eff}=H_0+[V_{-1},V_{+1}]/\hbar\Omega$ producing a static field in the
octupolar channel, and (iv) the paper's two key qualitative claims:
$h_m\propto\Gamma^{(3)}$ and anisotropy growth with drive strength $\zeta$.
\textbf{Verdict: REPLICATED (headline).}
\end{abstract}
\maketitle

\section{Target claim}
The headline (recipe \texttt{headline\_claim}) is that CPL induces an effective
static field linearly coupled to the magnetic octupole (OIFE) plus a
bond-dependent anisotropic exchange, opening a nonequilibrium multipolar phase
diagram. The quantitative kernel of the claim is the \emph{form, origin, and
magnitude} of the light-induced effective field $h_m$. We target this directly.

\section{Model reconstruction}
\subsection{$E_g$ non-Kramers doublet and pseudospin}
The $J=2$ manifold hosts the doublet (Eq.~1)
$\ket{\Uparrow}=(\ket{J_z{=}2}+\ket{J_z{=}{-}2})/\sqrt2$, $\ket{\Downarrow}=\ket{J_z{=}0}$.
Projecting the Stevens operators and normalizing (Eq.~2),
\begin{equation}
\frac{O_{22}}{4\sqrt3}\!\to\!\tilde\sigma^x,\quad
\frac{T_{xyz}}{2\sqrt3}\!\to\!\tilde\sigma^y,\quad
\frac{O_{20}}{12}\!\to\!\tilde\sigma^z,
\end{equation}
we recover \emph{exactly} the spin-$\tfrac12$ SU(2) algebra
$[\tilde\sigma^a,\tilde\sigma^b]=i\epsilon^{abc}\tilde\sigma^c$
(all three commutators verified to $<10^{-9}$; each projected operator equals
$\tfrac12$ Pauli). Crucially $\tilde\sigma^y$ is the \emph{magnetic octupole}
$T_{xyz}$ (time-reversal odd), while $\tilde\sigma^{x,z}$ are quadrupoles.

\subsection{Floquet effective couplings}
Using the paper's analytic forms (Eqs.~6a--6f), with $J_n$ the Bessel function
of the first kind, Floquet cutoff $p=7$, constraint $n{+}l{+}m=0$, and
representative parameters $t_{pd}{=}1.5$, $t_2{=}0.25$, $\tilde U{=}3.0$,
$\Delta_c{=}5.0$~eV, $\Omega\simeq0.414$~eV ($\sim100$~THz):
\begin{align}
J_\mathrm{eff}&=J^{(2)}-J^{(3)}+J^{(4)},\\
\Gamma^{(3)}&=\frac{t_{pd}^2 t_2}{9\sqrt3}\!\!\sum_{n+l+m=0}\!\!
\frac{\sin[(n{-}l)\psi_0]\,J_nJ_lJ_m}{(\tilde U{-}m\Omega)(\Delta_c{-}n\Omega)},\\
h_m&=\frac{t_{pd}^2 t_2}{8\sqrt3}\!\!\sum_{n+l+m=0}\!\!
\frac{\sin[(n{-}l)\psi_0]\,J_nJ_lJ_m}{(\tilde U{-}m\Omega)(\Delta_c{-}n\Omega)}.
\end{align}

\subsection{van Vleck mechanism}
On the doublet, a CPL drive
$V(t)=g[\tilde\sigma^x\cos\Omega t+\tilde\sigma^z\sin\Omega t]$ (two quadrupolar
channels, $90^\circ$ out of phase = circular) has Fourier components
$V_{\pm1}=\tfrac{g}{2}(\tilde\sigma^x\mp i\tilde\sigma^z)$. The high-frequency
correction $[V_{-1},V_{+1}]/\hbar\Omega\propto\tilde\sigma^y$ generates a static
field \emph{in the octupolar channel} --- the microscopic essence of the OIFE.

\section{Results}
\begin{table}[h]
\caption{Key verified quantities (full JSON in evidence/).}
\begin{ruledtabular}
\begin{tabular}{lc}
Quantity & Value\\\hline
SU(2) algebra $[\tilde\sigma^a,\tilde\sigma^b]{=}i\epsilon\tilde\sigma^c$ & confirmed ($<10^{-9}$)\\
$J_\mathrm{eff}$ at $\zeta{=}2$ & $1.47\times10^{-2}$ eV\\
$\Gamma^{(3)}$ at $\zeta{=}2$ & $-1.66\times10^{-4}$ eV\\
$h_m$ at $\zeta{=}2$ & $-1.87\times10^{-4}$ eV\\
$h_m/\Gamma^{(3)}$ (all $\zeta$) & $1.125=9/8$ (analytic)\\
$\Gamma^{(3)}/J_\mathrm{eff}$ growth $\zeta:0.5\!\to\!4$ & $7.7\!\times\!10^{-4}\!\to\!2.7\!\times\!10^{-3}$\\
$h_m$ small-$\zeta$ log--log slope & $1.96\approx2$ ($|E\times E^*|$)\\
van Vleck field channel & $\tilde\sigma^y$ (octupole)\\
\end{tabular}
\end{ruledtabular}
\end{table}

\begin{itemize}
\item \textbf{Form:} $h_m$ multiplies $\sum_i\tilde\sigma^y_i$ --- a uniform
field \emph{linear in the octupole}, exactly the OIFE structure of Eq.~(4).
\item \textbf{Origin:} $h_m$ vanishes at $\zeta{=}0$ and rises as $\zeta^2$
(fitted exponent $1.96$), i.e.\ $\propto|E(\Omega)\times E^*(\Omega)|$ ---
the optical-helicity / inverse-Faraday origin the paper asserts.
\item \textbf{Magnitude:} $h_m,\Gamma^{(3)}\sim10^{-4}$ eV and
$J_\mathrm{eff}\sim10^{-2}$ eV, consistent with the $10^{-2}$/$10^{-3}$ eV
scales of Fig.~3.
\item \textbf{Proportionality:} $h_m/\Gamma^{(3)}=9/8$ identically (from the
$1/8$ vs $1/9$ analytic prefactors), reproducing Fig.~3(b)'s statement that
$h_m$ and $\Gamma^{(3)}$ are proportional in this minimal model.
\end{itemize}

\section{Assessment}
The headline --- a CPL-induced effective \emph{static field linearly coupled to
the magnetic octupole}, of helicity ($\zeta^2$) origin, at the meV scale --- is
reproduced in form, origin, and magnitude from an independent build. Verdict
\textbf{REPLICATED (headline)}. Not attempted: the full many-body ED/DMRG phase
diagram (AFO/FO/PPFQ/IO/ML) and the absolute order-1 anisotropy of Fig.~3(a),
whose magnitude depends on geometry factors ($\psi_0$, $r_{pd}/r_{dd}$) not
fully specified in the excerpt (see \texttt{failure\_analysis.md}).

\section*{Kernel provenance}
Stevens operators and pseudospin construction reuse
\texttt{ollie\_multipolar\_stevens\_landau\_kernel.py} (Ollie's TEXTURES-100
multipolar kernel). Floquet coupling evaluation and van~Vleck demonstration are
original to this replication.

\end{document}
